% =====================================================================
%  Apresentação do modelo de estrutura
% =====================================================================
\chapter*[Apresentação do modelo]{Apresentação do modelo}
\addcontentsline{toc}{chapter}{Apresentação do modelo}

Este documento propõe uma nova forma de apresentar as ações do Instituto
Água e Terra (IAT) no Manual Operativo do Projeto (MOP) do PSH-PR. Toma-se
como piloto o Subcomponente 3.2.b -- \textit{Implantação e reabilitação de
sistemas coletivos de abastecimento de água e esgotamento sanitário no meio
rural} (item B.3.2.2 do MOP de 17 de março de 2026). O conteúdo técnico
foi preservado; o que muda é a organização. As atividades passam a ser
descritas como o capítulo de metodologia de uma dissertação: justificativa
e objetivos, área de estudo e critérios de seleção, delineamento, procedimentos
por atividade, salvaguardas, monitoramento e cronograma.

A mesma estrutura deve ser replicada nos demais subcomponentes sob
responsabilidade do IAT (Componente~1 e Subcomponente~3.2.b). O arquivo
\texttt{modelos/modelo-atividade.tex} contém o esqueleto em branco de uma
atividade.

\section*{Estrutura padrão do capítulo}

\begin{enumerate}[label=\arabic*.]
  \item \textbf{Contextualização e objetivos}: justificativa, objetivo geral,
        objetivos específicos e metas físicas.
  \item \textbf{Área de abrangência e critérios de seleção}: onde a ação
        acontece e como são escolhidos os territórios e beneficiários.
  \item \textbf{Delineamento metodológico}: premissas, encadeamento das
        atividades (fluxograma) e arranjo institucional e financeiro.
  \item \textbf{Procedimentos metodológicos por atividade}: uma seção por
        atividade, sempre com as mesmas subdivisões:
        \begin{enumerate}[label=\alph*)]
          \item objetivo e fundamentação;
          \item ficha-síntese (quadro padronizado);
          \item procedimentos, organizados em fases;
          \item etapas, produtos e evidências (quadro detalhado, quando houver);
          \item interfaces com outras atividades.
        \end{enumerate}
  \item \textbf{Salvaguardas ambientais e sociais}: riscos, mitigação e
        instrumentos do Quadro Ambiental e Social do Banco Mundial.
  \item \textbf{Monitoramento, avaliação e indicadores}: indicadores, metas,
        fontes de verificação e fluxo de validação.
  \item \textbf{Cronograma de execução}: diagrama de Gantt.
\end{enumerate}

\section*{Convenções de redação e formatação}

\begin{itemize}
  \item Redação impessoal, no futuro do presente (\textit{serão realizadas},
        \textit{será elaborado}), como em um projeto de pesquisa.
  \item Informação textual vai em \textbf{quadro} e dados numéricos vão em
        \textbf{tabela} (ABNT NBR 14724). Todo quadro, tabela e figura traz
        a fonte.
  \item Os códigos das atividades (3.2.2.1 a 3.2.2.6) seguem o MOP e são
        citados com o comando \verb|\atividade{3.2.2.x}|.
  \item Inconsistências encontradas durante a reestruturação são marcadas
        com \verb|\pendencia{...}| e listadas no Apêndice~\ref{ap:pendencias}.
        Para ocultá-las na versão final, basta desativar a opção em
        \texttt{main.tex}.
\end{itemize}

\section*{Correspondência com o MOP vigente}

O \autoref{qua:correspondencia} indica onde o conteúdo do item B.3.2.2 do MOP
foi reposicionado na nova estrutura.

\begin{quadro}[htbp]
\caption{Correspondência entre o MOP (versão de 17/03/2026) e a estrutura proposta}
\label{qua:correspondencia}
\small
\renewcommand{\arraystretch}{1.2}
\begin{tabularx}{\textwidth}{|L|L|}
\hline
\cab{MOP -- item B.3.2.2} & \cab{Estrutura proposta} \\ \hline
Texto introdutório e objetivos específicos & \ref{sec:contexto} Contextualização e objetivos \\ \hline
Frentes de trabalho (83 comunidades; 2.300 sistemas) & \ref{subsec:metas} Metas \\ \hline
Área de abrangência e critérios de seleção & \ref{sec:area} Área de abrangência e critérios de seleção \\ \hline
Executor, forma de execução, prazo e custo & \ref{subsec:arranjo} Arranjo institucional e financiamento \\ \hline
Quadro 26 -- Aspectos socioambientais & \ref{sec:salvaguardas} Salvaguardas ambientais e sociais \\ \hline
Quadro 27 -- Indicadores & \ref{sec:monitoramento} Monitoramento, avaliação e indicadores \\ \hline
B.3.2.2.1 a B.3.2.2.6 (atividades) & \ref{sec:procedimentos} Procedimentos metodológicos por atividade \\ \hline
Quadros 28 e 29 -- Etapas das atividades & Quadros de etapas dentro de cada atividade \\ \hline
Quadro 30 -- Atividades, responsáveis e prazos & Fichas-síntese de cada atividade e \ref{sec:cronograma} Cronograma \\ \hline
Quadro 31 -- Produtos e evidências & Fichas-síntese de cada atividade \\ \hline
\end{tabularx}
\fonte{Elaborado pelo IAT (2026).}
\end{quadro}
